\documentclass[aspectratio=169,10pt,t]{beamer}
\usepackage{ac1}
\aulainfo{29}{Pandas na Prática}{Agrupamentos, agregações, valores ausentes e gráficos}

\begin{document}
\capa

\begin{frame}{Roteiro}
  \begin{itemize}
    \item Cálculos agrupados: \texttt{groupby} e agregações
    \item Contagem de frequências: \texttt{value\_counts}, \texttt{nunique}
    \item Valores ausentes: \texttt{isna}, \texttt{fillna}, \texttt{dropna}
    \item Datas e horas
    \item Gráficos direto do Pandas
    \item Roteiro de uma análise exploratória
  \end{itemize}
\end{frame}

\begin{frame}{Perguntas típicas de quem analisa dados}
  \begin{itemize}
    \item Qual o tempo de ciclo \textbf{médio de cada máquina}?
    \item Qual a \textbf{taxa de peças com defeito} por turno?
    \item Quantas peças \textbf{de cada tipo} cada operador produziu?
    \item (Gapminder) Qual foi a expectativa de vida média \textbf{por ano}? E \textbf{por continente}?
  \end{itemize}
  \begin{block}{Padrão: dividir--aplicar--combinar (\emph{split--apply--combine})}
    \textbf{1.} dividir os dados em grupos \(\to\) \textbf{2.} aplicar um cálculo a cada grupo \(\to\) \textbf{3.} combinar os resultados numa tabela.
  \end{block}
\end{frame}

\begin{frame}[fragile]{\texttt{groupby}: médias agrupadas}
\begin{lstlisting}
media_por_maquina = df.groupby("maquina")["tempo_ciclo_seg"].mean()
print(media_por_maquina.round(2))
\end{lstlisting}
\begin{lstlisting}[style=saida]
maquina
M1    32.92
M2    28.34
M3    45.28
M4    30.25
Name: tempo_ciclo_seg, dtype: float64
\end{lstlisting}
  \begin{itemize}\small
    \item \texttt{groupby("maquina")} agrupa; \texttt{["tempo\_ciclo\_seg"]} escolhe a coluna; \texttt{.mean()} calcula
    \item O objeto agrupado sozinho não mostra nada útil: falta aplicar a agregação
  \end{itemize}
\end{frame}

\begin{frame}[fragile]{Várias estatísticas de uma vez: \texttt{agg}}
\begin{lstlisting}[basicstyle=\ttfamily\scriptsize]
resumo = (
    df
    .groupby("maquina")["tempo_ciclo_seg"]
    .agg(["count", "mean", "std", "min", "max"])
    .round(2)
)
print(resumo)
\end{lstlisting}
\begin{lstlisting}[basicstyle=\ttfamily\scriptsize,style=saida]
         count   mean   std   min   max
maquina
M1          17  32.92  3.55  28.0  38.9
M2          18  28.34  3.49  20.9  33.6
M3          12  45.28  3.84  39.9  52.4
M4          13  30.25  2.00  26.8  34.6
\end{lstlisting}
  \small Os parênteses externos permitem o \emph{encadeamento de métodos} em várias linhas (estilo preferido, como no notebook).
\end{frame}

\begin{frame}[fragile]{Agrupando por mais de uma coluna}
\begin{lstlisting}
tab = df.groupby(["maquina", "status"]).size().unstack(fill_value=0)
print(tab)
\end{lstlisting}
\begin{lstlisting}[style=saida]
status   Aprovado  Descarte  Retrabalho
maquina
M1             13         1           3
M2             13         0           5
M3             10         0           2
M4             12         0           1
\end{lstlisting}
  \begin{itemize}\small
    \item Agrupar por duas colunas gera um \emph{índice composto}; \texttt{unstack()} transforma o segundo nível em colunas
    \item \texttt{reset\_index()} volta a ter um DataFrame ``plano'' (os nomes viram colunas comuns)
    \item No notebook: \texttt{df.groupby(["year", "continent"])[["lifeExp", "gdpPercap"]].mean()}
  \end{itemize}
\end{frame}

\begin{frame}[fragile]{Taxa de defeitos por máquina e por turno}
\begin{lstlisting}
df["com_defeito"] = df["status"] != "Aprovado"          # True/False
print(df.groupby("maquina")["com_defeito"].mean().round(2))
print(df.groupby("turno")["com_defeito"].mean().round(3))
\end{lstlisting}
\begin{lstlisting}[style=saida,basicstyle=\ttfamily\scriptsize]
maquina: M1 0.24 | M2 0.28 | M3 0.17 | M4 0.08
turno:   Manha 0.222 | Noite 0.240 | Tarde 0.118
\end{lstlisting}
  \begin{itemize}\small
    \item A média de uma coluna booleana (\(True=1\)) é a \alert{proporção} de \texttt{True}
    \item M2 tem a maior taxa de problemas, apesar de M3 ser a mais lenta: tempo e qualidade são perguntas diferentes!
  \end{itemize}
\end{frame}

\begin{frame}[fragile]{Frequências: \texttt{value\_counts} e \texttt{nunique}}
\begin{lstlisting}
print(df["maquina"].value_counts())            # M2 18, M1 17, M4 13, M3 12
print(df["operador"].nunique())                # quantos operadores distintos
print(df.groupby("turno")["operador"].nunique())
print(df["defeito"].value_counts())            # ignora NaN por padrão
print(df["turno"].value_counts(normalize=True).round(2))   # proporções
\end{lstlisting}
  \begin{itemize}
    \item \texttt{value\_counts()}: frequência de cada valor (ordenada)
    \item \texttt{nunique()}: número de valores \emph{distintos}
    \item No \emph{gapminder}: \texttt{df.groupby("continent")["country"].nunique()} \(\to\) África 52, Américas 25, Ásia 33, Europa 30, Oceania 2
  \end{itemize}
\end{frame}

\begin{frame}[fragile]{Valores ausentes}
\begin{lstlisting}
print(df.isna().sum())               # ausentes por coluna
# defeito: 48 ausentes  (= as 48 peças aprovadas)

df["defeito"] = df["defeito"].fillna("sem defeito")   # preenche
df.dropna(subset=["tempo_ciclo_seg"])                 # remove linhas
print(df["defeito"].value_counts().head(3))
\end{lstlisting}
\begin{lstlisting}[style=saida]
defeito
sem defeito    48
Risco           5
Trinca          4
\end{lstlisting}
  \begin{alertblock}{Decida com critério}
    Ausente significa ``não se aplica'' (aqui: peça sem defeito) ou ``não foi medido''? Preencher, remover ou manter muda os resultados: documente a decisão.
  \end{alertblock}
\end{frame}

\begin{frame}[fragile]{Datas e horas}
\begin{lstlisting}
df["timestamp"] = pd.to_datetime(df["timestamp"])
print(df["timestamp"].min(), df["timestamp"].max())
# 2026-08-03 06:04:00   2026-08-03 12:54:00
df["hora"] = df["timestamp"].dt.hour
pecas_por_hora = df.groupby("hora").size()
print(df["timestamp"].dt.day_name().head(1))    # Monday
\end{lstlisting}
  \begin{itemize}
    \item \texttt{pd.to\_datetime} converte texto em data/hora; o acessor \texttt{.dt} extrai \texttt{year}, \texttt{month}, \texttt{day}, \texttt{hour}\ldots
    \item Permite agrupar por mês, semana ou hora, e calcular \emph{intervalos} entre eventos
    \item Em dados reais, datas costumam vir em formatos variados: confira \texttt{dtypes}
  \end{itemize}
\end{frame}

\begin{frame}[fragile]{Gráficos direto do Pandas}
  \begin{columns}[T]
    \column{0.5\textwidth}
\begin{lstlisting}[basicstyle=\ttfamily\scriptsize]
import matplotlib.pyplot as plt

# média por máquina (barras)
(df.groupby("maquina")["tempo_ciclo_seg"]
   .mean()
   .plot(kind="bar", rot=0))
plt.ylabel("Tempo médio (s)")
plt.show()

# histograma de uma coluna
df["tempo_ciclo_seg"].plot(kind="hist",
                           bins=10)
\end{lstlisting}
    \column{0.47\textwidth}
    \includegraphics[width=\textwidth]{fig_media.png}
  \end{columns}
  \begin{itemize}\small
    \item \texttt{.plot(kind="bar" | "hist" | "box" | "line" | "scatter")}: usa o Matplotlib por baixo
    \item Para personalizar, use \texttt{fig, ax = plt.subplots()} e passe \texttt{ax=ax} ao \texttt{.plot}
  \end{itemize}
\end{frame}

\begin{frame}[fragile]{Série temporal do \emph{gapminder}}
\begin{lstlisting}
expect_global = df.groupby("year")["lifeExp"].mean()
print(expect_global.head(3))
# year
# 1952    49.057620
# 1957    51.507401
# 1962    53.609249

expect_global.plot()
plt.ylabel("Expectativa de vida média (anos)")
plt.show()
\end{lstlisting}
  \begin{itemize}
    \item Exemplo final do notebook: o gráfico de linha mostra a \textbf{tendência de crescimento} (de 49 para 67 anos entre 1952 e 2007)
    \item Visualização e agrupamento \alert{juntos}: é assim que se explora dados
  \end{itemize}
\end{frame}

\begin{frame}{Roteiro de uma análise exploratória}
  \begin{enumerate}
    \item \textbf{Carregar} (\texttt{read\_csv}) e \textbf{conhecer}: \texttt{shape}, \texttt{dtypes}, \texttt{head}, \texttt{describe}
    \item \textbf{Limpar}: ausentes, tipos (\texttt{to\_datetime}), duplicatas, valores impossíveis
    \item \textbf{Perguntar}: o que eu quero saber? (taxa, média, comparação, tendência)
    \item \textbf{Calcular}: filtros, \texttt{groupby}, \texttt{agg}
    \item \textbf{Visualizar}: o gráfico certo para cada pergunta (Aula 27)
    \item \textbf{Concluir e documentar}: o que os dados mostram? Quais as limitações?
  \end{enumerate}
  \begin{exampleblock}{Este é o esqueleto de muitos projetos finais}
    Combine: arquivos CSV \(+\) Pandas \(+\) funções bem organizadas \(+\) gráficos \(+\) testes (Aulas 30 a 33).
  \end{exampleblock}
\end{frame}

\begin{frame}{Atividade no laboratório}
  Com \texttt{dadosproducao.csv}:
  \begin{enumerate}
    \item Calcule o tempo médio por máquina \emph{e} por turno (tabela com as duas chaves).
    \item Qual operador produziu mais peças com problema (\texttt{Retrabalho} ou \texttt{Descarte})?
    \item Qual o defeito mais comum em cada máquina?
    \item Quantas peças foram produzidas por hora? Faça um gráfico de linha.
    \item Escreva uma função \texttt{relatorio(df)} que imprime um resumo (totais, taxas e a pior máquina) e uma função que salva dois gráficos em PNG.
  \end{enumerate}
  \footnotesize Dica: para o item 3 use \texttt{df.groupby(["maquina","defeito"]).size()}.
\end{frame}

\resumo{
  \item \texttt{groupby} + agregação responde perguntas por grupo; \texttt{agg} calcula várias estatísticas
  \item \texttt{value\_counts} e \texttt{nunique} fazem contagens de frequências
  \item Valores ausentes: \texttt{isna}, \texttt{fillna}, \texttt{dropna}; decida e documente
  \item \texttt{.plot()} cria gráficos direto de Series e DataFrames
  \item Próxima aula: especificação do projeto final
}

\end{document}
